\chapter{Conclusion}
We implemented an optimized version of the Snowman application which uses tile decomposition and dynamic scheduling. We could show that our implementation performs better than the unoptimized code when keeping the following considerations in mind:\\
\\
Using dynamic scheduling introduces extra overhead, since there is a master process that does not contribute to actually rendering the image. Therefore, dynamic scheduling does not perform well when using very few processes. Our experiments showed that using eight or more processes with dynamic scheduling led to better performance than the unoptimized implementation. Another important thing to consider is the size of the tiles. It needs to be chosen so that there are enough tiles to keep all workers busy. Our results suggest than when using 96 processes, one should use a tile width of 64 pixels or less. When using less processes, larger tile sizes may be acceptable.\\
\\
Notably, our implementation does not have the "spiky" runtime fluctuations that we saw in the original implementation. This is because of better handling of the remainder when the image size can not be evenly divided by the number of workers. Another problem with the original implementation was that some workers get assigned an entirely white section of the image, while others get sections that are more complex to render. Tile decomposition avoids this issue, which presumably also contributes to some of the speedup we saw.